% Generated draft: requires mathematical and source review.
\documentclass[11pt]{article}
\usepackage[margin=25mm]{geometry}
\usepackage{fontspec}
\setmainfont{Latin Modern Roman}
\usepackage{amsmath,amssymb,mathtools,unicode-math}
\setmathfont{Latin Modern Math}
\usepackage{enumitem,needspace,xurl,hyperref}
\setlength{\parindent}{0pt}
\setlength{\parskip}{5pt}
\setlength{\emergencystretch}{3em}
\newcommand{\N}{\mathbb{N}}
\newcommand{\Z}{\mathbb{Z}}
\newcommand{\Q}{\mathbb{Q}}
\newcommand{\R}{\mathbb{R}}
\newcommand{\C}{\mathbb{C}}
\newcommand{\F}{\mathbb{F}}
\newcommand{\A}{\mathbb{A}}
\newcommand{\PP}{\mathbb{P}}
\newcommand{\E}{\mathbb{E}}
\newcommand{\eps}{\varepsilon}
\newcommand{\dd}{\,\mathrm{d}}
\begin{document}
\section*{UnsolvedMath ID: 20002710}
\textbf{Source problem code:} catalogue entry
\subsection{Definitions and mathematical statement}
H. T. Hall's problem, extracted with the split label ``Problem 2. 4.,'' starts from classical meanders across $2n$ ordered bridges. Two noncrossing pairings of $[2n]$, one on each side, form a meandric system; a meander has one connected component. The source recalls a bijection with ordered pairs of type-$A_{n-1}$ noncrossing partitions at maximal distance in the undirected Hasse graph, and asks whether it yields new counting methods. It proposes type-$W$ meanders as ordered diameter pairs in $NC_W$, asks for a combinatorial model, and specifically asks about type $B$.

For a finite Coxeter system of rank $r$, fix a Coxeter element $c$, let $NC_c(W)=[1,c]$ in absolute order, and let $d_H$ be distance in its undirected Hasse graph. Its diameter is $r$; an ordered diameter is $(u,v)$ with $d_H(u,v)=r$, with endpoint order counted. The aid's type-$B$ result is: for every $n\geq3$, exactly
\[
n^3(n+1)
\]
ordered diameters in $NC_c(B_n)$ have endpoint rank pairs $(1,n-1)$ or $(n-1,1)$, counted together. Adding $(1,c)$ and $(c,1)$ gives the lower bound $n^3(n+1)+2$ for all ordered diameters. For $B_3$, these are all possibilities, giving exactly $110$. For larger $n$ the statement does not count additional middle-rank diameter pairs.

The aid also gives a $B_3$ combinatorial interpretation: view a type-$B$ noncrossing partition as a centrally symmetric type-$A$ partition on the signed points, apply the standard doubling/fattening to obtain noncrossing pairings, and superpose the pairings for an ordered pair. The resulting doubled meandric system is centrally symmetric. The $B_3$ diameter pairs correspond precisely to those systems whose every connected component is fixed setwise by the half-turn. This permits several components and must not be replaced by the classical one-component condition. The supplied result is stated for $B_3$, not as a complete general-$B_n$ characterization.

The source remarks, attributed to A. Nica and J. Scott, ask about a type-$W$ meander determinant and its factorization, and about links to chromatic polynomials and Temperley--Lieb algebras. For an explicit distance convention consistent with the aid's formula, use
\[
M_W(q)_{u,v}=q^{r+1-d_H(u,v)}.
\]
Different overall powers would change determinant factors. For the dihedral group $I_2(m)$, $m\geq3$, rank is two, the poset has $m$ middle elements and two bounds, and the aid's exact results are
\[
\#\{\text{ordered diameters}\}=m(m-1)+2,
\]
\[
\det M_{I_2(m)}(q)
=q^{m+2}(q^2-1)^{m+1}(q^2-(m-1)).
\]
Here $q$ is an indeterminate and the matrix entries are $q^{3-d_H}$. The irreducible dihedral range is explicit; extending the notation to the reducible $m=2$ case is a separate convention. The source's type-$A$ Chebyshev factorization and these dihedral formulas do not give a complete type-$W$ factorization or settle the chromatic/Temperley--Lieb significance question. The trailing ``3. Reflection Groups'' is a new section heading, not another assumption.
\subsection{Short English statement}
The source asks for Coxeter-type meander counts, models, and determinant factorizations. The aid counts all extreme-rank type-$B$ diameters, obtains the complete $B_3$ count and a symmetry model there, and factors the distance-meander determinant in every irreducible dihedral type. General counting, middle-rank type-$B$ behavior, and the broader algebraic links remain separate questions.
\subsection{Sources}
\begin{itemize}
\item Local UnsolvedMath registry, permanent ID $20002710$, \texttt{catalogue entry}; \texttt{aim-probability-notes.json}, \texttt{notes[151]}, extracted source item 2.
\item American Institute of Mathematics, January 2005 \emph{Braid Groups, Clusters, and Free Probability}, H. T. Hall meander problem; supplied verification reports pages 6--7, \url{https://aimath.org/WWN/braidgroups/braidgroups.pdf}.
\item Original remarks by A. Nica and J. Scott: type-$A$ meander determinants factor into Chebyshev polynomials, reference [20]; chromatic-polynomial and Temperley--Lieb connections, reference [15]. Both bibliography entries are absent from the supplied record.
\item Research attempt 1: partial result, problem partially solved at run, rigor label ``proved,'' novelty confidence low. Registry classification partial progress, status partially solved, statement status exact. No authoritative separate citation or explicit enumeration certificate is supplied for the type-$B$ and dihedral results.
\end{itemize}
\end{document}
